\documentclass[11pt]{article}
\usepackage[a4paper,margin=27mm]{geometry}
\usepackage{amsmath,amssymb,amsthm,hyperref}
\newtheorem{theorem}{Theorem}
\newtheorem{lemma}{Lemma}
\newtheorem{corollary}{Corollary}
\title{A square-root obstruction to uniform endpoint-profile rates}
\author{Research proof; independent mathematical review completed}
\date{30 September 2026}
\begin{document}\maketitle
\noindent\textbf{Scope.} We first give a counterexample in the general
finite, positively weighted monotone comparison-tensor class, then
transfer its robust inequality to completely finite equiprobable dice.
It concerns the rate of agreement of individual columns, rather than
the already sharp additive offset in the one-sided relative bound.

\begin{theorem}\label{main}
For all sufficiently large $m$ there is a finite monotone tensor
$p_i(q)\in[0,1]$, $1\le i\le m$, with positive row weights $w_q$
summing to one and
\[
 \sum_qw_q\prod_{i\in S}p_i(q)=\frac1{|S|+1}
                      \qquad(S\subseteq[m]),
\]
having a row with $1\le s:=\sum_i(1-p_i)\le100$ at which
\[
 \max_i|m(1-p_i)-s|\ge c m^{-1/2}.
\]
All but one of the columns are identical. Thus a universal endpoint
profile estimate $o(m^{-1/2})$ on bounded $s$ is impossible in this
tensor class.
\end{theorem}

\begin{lemma}\label{gauss}
Let $Q=\lceil(m+1)/2\rceil$. The $Q$-node Gauss--Legendre probability
rule on $[0,1]$ has a node $a$ and weight $v$ with
\[
 m(1-a)\longrightarrow\pi^2,\qquad v\ge c m^{-3/2}.
\]
\end{lemma}
\begin{proof}
Number the nodes downwards from $1$, and choose index
$j=\lfloor\sqrt m\rfloor$. Write $2a-1=\cos\theta$.
The elementary zero bounds in
\href{https://dlmf.nist.gov/18.16.E2}{DLMF 18.16.2}, specialized
at Jacobi parameters $0,0$, give
\[
 \frac{(j-1/2)\pi}{Q+1/2}\le\theta
       \le\frac{j\pi}{Q+1/2}.
\]
Thus $\theta\sim2\pi/\sqrt m$.
Since $1-a=(1-\cos\theta)/2$, the
asserted location follows.

For the weight, no quadrature-weight asymptotic is needed.
The Bernstein inequality for Legendre polynomials is
\[
 (1-x^2)^{1/4}|P_r(x)|
 \le \sqrt{2/\pi}\,(r+1/2)^{-1/2}\quad(-1<x<1).
\]
This is \href{https://dlmf.nist.gov/18.14.E7}{DLMF 18.14.7}
with Gegenbauer parameter $1/2$ (also 18.14.3 at Jacobi parameters
$0,0$). Hence the shifted Legendre reproducing kernel satisfies
\[
 K_{Q-1}(a,a)=\sum_{r=0}^{Q-1}(2r+1)P_r(2a-1)^2
                  \le\frac{C Q}{\sin\theta}.
\]
At a Gaussian node its probability weight is
$v=K_{Q-1}(a,a)^{-1}$. To verify the identity, the kernel polynomial
at that node, divided by its diagonal value, is the cardinal
Lagrange polynomial there; integrate its square using exactness
through degree $2Q-2$ and the reproducing identity.
Thus $v\ge c\sin\theta/Q\ge c m^{-3/2}$.
\end{proof}

\begin{proof}[Proof of Theorem~\ref{main}]
Let $\sigma=\sum_{\ell=1}^Qv_\ell\delta_{a_\ell}$ be the Gaussian rule, and $F$ its
cumulative distribution function. It integrates every polynomial of
degree at most $m$, because $2Q-1\ge m$.
On the continuous row parameter $t\in[0,1]$, with uniform row measure,
set
\[
 p_1(t)=F(t),\qquad p_2(t)=\cdots=p_m(t)=t.
\]
All columns are nondecreasing. A mixed moment not involving column
$1$ is immediate. If it involves column $1$ and $r$ other columns,
$0\le r\le m-1$, then Fubini gives
\[
 \int_0^1F(t)t^r\,dt
 =\frac{1-\int a^{r+1}\,d\sigma(a)}{r+1}
 =\frac1{r+2}.
\]
Thus all required mixed moments hold.

At the node $a$ of Lemma~\ref{gauss}, the two one-sided values of
$F$ differ by $v$. At least one, denoted $F_*$, satisfies
$|F_*-a|\ge v/2$. The corresponding comparison row has
\[
 s=m(1-a)+(a-F_*),\qquad
 m(1-p_1)-s=(m-1)(a-F_*).
\]
The first lies in $[1,100]$ for large $m$, since $|a-F_*|\le1$;
the second has absolute value at least $c m^{-1/2}$.

Finally make the row measure finite without losing this row.
Partition $[0,1]$ at all nodes $a_\ell$. On each resulting interval,
$F$ is constant, so every mixed-moment integrand above is a polynomial
of degree at most $m$. Apply a positive Gauss--Lobatto quadrature
of degree at least $m$ separately on every interval, retaining both
endpoints. At a shared endpoint keep the two resulting rows distinct,
with the left and right values of $F$, respectively. Their other
$m-1$ coordinates agree, and ordering the lower $F$ value first
preserves monotonicity. The positive weights sum to one and preserve
every mixed moment. Both one-sided comparison rows at $a$ occur with
positive weight, so the required one is retained.
\end{proof}

\begin{corollary}[Actual finite uniform dice]\label{dice}
For all sufficiently large $m$ there is a permutation-fair family of
$m+1$ finite uniform dice and a face of one die whose comparison
probabilities satisfy
\[
 1\le s=\sum_{i=1}^m(1-p_i)\le100,\qquad
 \max_i|m(1-p_i)-s|\ge c m^{-1/2}.
\]
No bound on the face counts is asserted here.
\end{corollary}
\begin{proof}
Use the Gaussian measure $\sigma$, node $a$, weight $v$, and
favorable one-sided CDF value $F_*$ in the preceding proof.
Choose $t\in(0,1)$ sufficiently close to $a$ on that side, and
strictly between Gaussian nodes, so that $|t-a|<v/8$ and
$F_\sigma(t)=F_*$. Also require $m|t-a|<1/10$; this is possible
by arbitrarily close approach. Then
\[
 |F_\sigma(t)-t|\ge3v/8,
 \qquad m(1-t)\longrightarrow\pi^2.
\]
(For the displayed limit one may choose $m|t-a|\to0$.)
Apply the realization theorem in
\path{finite_profile_realization.tex} with $\rho<v/100$ and,
if necessary, still smaller so that $m\rho<1/10$.
For the resulting comparison row, put $e_i=p_i-t$ for $i\ge2$.
Then
\[
 m(1-p_1)-s=(m-1)(t-p_1)+\sum_{i=2}^m e_i.
\]
Its absolute value is at least
$(m-1)(3v/8-2\rho)\ge c m v\ge c m^{-1/2}$.
Moreover
\[
 s=m(1-t)+(t-p_1)-\sum_{i=2}^m e_i
\]
lies in $[1,100]$ for large $m$, since $|t-p_1|\le1$.
Exact permutation fairness and uniform face probabilities come from
the realization theorem.
\end{proof}

\paragraph{Matching upper bound.}
The subsequently proved theorem in
\path{independent_directions/natural_scale_endpoint_profile.tex}
gives
\[
 \max_i|m(1-p_i)-s|\le C_L\bigl(\sqrt{s/m}+m^{-1}\bigr)
 \qquad(0\le s\le L)
\]
for each fixed $L$, including arbitrary positive row weights.
The present examples force the square-root order at fixed positive $s$;
\path{endpoint_profile_sharpness.tex} separately forces the inverse-$m$
additive floor, even for actual finite uniform dice. The subsequent note
\path{independent_directions/intermediate_endpoint_sharpness.tex}
constructs matching finite examples at every bounded positive endpoint
scale, with realized deficit comparable to the prescribed scale.
\end{document}
